\scanpage{4}
% Broad oblique strokes cancel the upper-left working, including its diagram.
\[
\begin{aligned}
 &\deleted{\unsure{\mu_p|\mathbb{H}_p}}\quad \deleted{G}\quad \deleted{G^0}\qquad \deleted{\unsure{k[[t|k[[+}}\\
 &\hspace{2em}\deleted{k}\qquad \deleted{\deleted{\unsure{\mathcal O|G}}\unsure{k_{\mathrm{red}}|G_{\mathrm{red}}}}
 \hspace{9em}\deleted{\Delta_0\unsure{\ni|\supset|\Rightarrow}\Delta}
\end{aligned}
\]

\[
\begin{aligned}
 \deleted{\unsure{\mathcal R|R}=\unsure{\mathbb A b|\mathrm{Ab}}(\unsure{\mathcal C|C})^{\mathrm{op}}}
 &\quad\deleted{\unsure{\overline{\mathcal R}|\overline{R}}}
 \qquad\deleted{\sigma=\{G^{\unsure{n|h}}\}}\\
 \deleted{\unsure{\prime\prime|\prime}}\quad &\deleted{G^n\unsure{-|\longrightarrow}G^m}\\
 \deleted{\mathbb A b(\mathcal C)}\quad &\deleted{G^n\longrightarrow G}
\end{aligned}
\]

\diagram{\begin{tikzcd}[ampersand replacement=\&,row sep=small,execute at end picture={\draw[gray] (\tikzcdmatrixname-3-1.south west) -- (\tikzcdmatrixname-1-1.north east);\draw[gray] (\tikzcdmatrixname-3-2.south west) -- (\tikzcdmatrixname-1-3.north east);}]
{\deleted{\text{\uuline{$\mathrm{Hom}$}}'(\text{\textcircled{$\sigma$}},(\mathrm{ENS})\unsure{)|{}}}}\arrow[dd,bend left=5] \& \& {\deleted{\unsure{FN|EN}}}\\
{} \& {\deleted{\text{\unsure{topos|transp.}}}}\arrow[l] \& {}\\
{\deleted{(\mathrm{ENS})}}\arrow[uu,bend left=20] \& {\deleted{\unsure{\mathrm{sp}|\mathrm{sr}}\;\mathcal R\Rightarrow}} \& {}
\end{tikzcd}}

% The following right-hand block has three diagonal cancellation strokes.
\deleted{\deleted{Soit} Disons posant $(C'_i,\lambda'_i)=(C_{\psi(i)},\lambda_{\psi(i)})$, on voit que \deleted{$C'_i$ est tel que}}

\diagram{\begin{tikzcd}[ampersand replacement=\&,row sep=small,execute at end picture={\draw[gray] (\tikzcdmatrixname-2-1.south west) -- (\tikzcdmatrixname-1-2.north west);}]
{\deleted{C'_{i\deleted{\unsure{+1|+n}}}}}\arrow[r] \& {\deleted{C'_{i+1}}}\\
{\deleted{C_{\psi(i)}}}\arrow[u,equal] \& {\deleted{C_{\psi(i+1)}=g\psi(i)}}\arrow[u,equal]
\end{tikzcd}}

\deleted{est $\lambda'_i$-préexact. Donc si $\omega'$ est le plus petit ordinal $<\omega$ tel que $\omega'$ soit de \unsure{cardinalité|cofinalité} $\pi'$, on voit que \deleted{\unsure{si}} (par le \inserted{\unsure{=}} lemme du 1°) $\lambda_{\omega}^{\prime\unsure{\star|{}}}$ est exact dans $C'_{\omega}$, \deleted{donc $g(\omega)=$} cqfd.}

\margin{\uuline{\unsure{Lemme}}}
\uline{Dém} prop \qquad a) est trivial \qquad\unsure{$\sim$|v}

b) \deleted{Soit $x\in$} OPS $U,V\ne\varnothing$, \deleted{\illegible{1} si pour} on est ramené à prouver que si \deleted{$x\in W=U.V\quad W(k)$} $x=u_0v_0$ \inserted{$(u_0\in U(k),v_0\in V(k))$} \deleted{(tel à $U.V$)} alors l'image inverse de $\unsure{UV|U.V}$ dans $G$ par $G\xrightarrow{\varphi}X$\quad $\varphi(g)=g.x$, \deleted{est st} contient $G_0$. i.e. que pour $g\in G^0$, \unsure{\ldots} $gx=u.v$ (après ext. de la base \deleted{p. $u\in G^0(k')$, $v\in$} $u\in U(k')$, $v\in V(k')$)

\deleted{\unsure{Ou|Or}} Ou encore $u^{-1}gx\in V$ i.e. $u^{-1}g\in\varphi^{-1}(W)=W$ i.e. \deleted{$u\in g$} $u\in gW^{-1}$, donc il suffit de voir que $gW^{-1}\cap U\ne\varnothing$. Or soit $T=G^0u_0$, alors \deleted{$U\cap$} $T$ est irréductible et, \deleted{$T\cap U\ne\varnothing$} $T\cap U\ni u_0$, $gW^{-1}\cap\deleted{U}T\ni gu_0$ [car $gu_0\in T$, $gu_0\in gW^{-1}$ i.e. $u_0\in W^{-1}$ ou $u_0^{-1}\in W$ car $u_0^{-1}.x=v_0\in V$].

Mais alors $gW^{-1}\cap T$ et $U\cap T$ sont des \unsure{ouverts|ouvert} non vides de $T$ irréductible, donc \deleted{\unsure{sont intersectés|se coupent}} \inserted{\unsure{int.|inter.} $\ne\varnothing$, OK}

c) Résulte de la partie a) du

\uline{Lemme} \deleted{Soit \unsure{T|J} une partie de $X$ stable par $G^0$ \unsure{est} stable par spécialisation.} \deleted{\unsure{a)}}\inserted{Si les $\varphi_x:G_{\unsure{\kappa|x}}\longrightarrow X_{\unsure{\kappa|x}}$ \unsure{gén.}} \deleted{Si $\pi$ générisant} (p. ex. $X$ un ps torseur) alors une partie $W$ de $X$ est stable par $G^0$ ssi elle est stable par spécialisation et par générisation (i.e., si $X$ loc \unsure{noeth|nœth}, ssi elle est ouverte et fermée)

\deleted{\unsure{Dém|Réc}} \unsure{b)|1)} Supposons $W$ stable par $G^0$, \inserted{i.e. $\pi$ induit $G^0\times W\longrightarrow W$, donc} \deleted{soit \unsure{$x\in$|$x$} $x,y\in X$ \unsure{tq}, $y\in\overline{\{x\}}$, prouvons} \deleted{i.e. $G\times X\longrightarrow X\times X$ induit $G^0\times W\longrightarrow W\times W$.} \deleted{que $x\in W\Longleftrightarrow y\in W$. \unsure{1)} $x\in W$ prouvons $y\in W$. OPS $x\in W(k)$,} \margin{\deleted{Or $G\times W\to W\times W$ est générisant et $G$}} \deleted{Mais alors Soit en effet} \deleted{Si $W$ est st. Mais $G^0\times W$ est stable}

Montrons que si $y\in W$, $x$ générise $y$, alors $x\in W$. Or OPS $y\in W(k)$, et on a alors $\varphi:g\longmapsto gy$\quad $G\longrightarrow X$, induisant $\unsure{\varphi'|g'}:G^0\longrightarrow W$, or $\varphi$ est \unsure{générisant|generisant} et $G^0$ stable par générisation, donc $\varphi(G^0)$ stable par générisation. On \unsure{a|.} prouvé que $W$ est stable par générisation, donc $\complement W$ l'est aussi, donc $W$ est stable par spécialisation
% archive sheet number: 3

